%
% Chapter 2 — Background and Related Work
%
\chapter{Background} \label{cap:background}

\section{Digital Journalism Platforms} \label{sec:background_journalism}

Digital journalism has undergone significant transformations in the past decade, driven
by the proliferation of social media, mobile devices, and algorithmic content
distribution. Traditional newsrooms have increasingly migrated to web-native platforms,
yet many existing solutions prioritise speed and volume over editorial rigour and
journalistic depth~\cite{slow_journalism}.

Slow journalism, as a counter-movement, emphasises long-form reporting, investigative
depth, and contextual analysis. Platforms designed to support slow journalism require
features that differ substantially from high-volume news aggregators: robust editorial
workflows, multimedia support, and fine-grained access control for editorial teams.

\section{Existing Solutions} \label{sec:background_existing}

Several content management systems (CMS) were analysed as potential off-the-shelf
solutions before the decision was made to develop a bespoke platform.

\subsection{WordPress} \label{sec:background_wordpress}

WordPress is the most widely adopted CMS in the world, powering over 40\% of all
websites~\cite{wordpress_stats}. A WordPress-based proof of concept was the starting
point for this project. Although WordPress provides a rich plugin ecosystem and a
familiar editorial interface, it presents notable limitations in the context of
Diário LX:

\begin{itemize}
      \item Limited control over the editorial approval workflow without heavy plugin
            dependencies.
      \item Performance and security concerns at scale.
      \item Architectural rigidity that complicates custom API design and frontend
            decoupling.
\end{itemize}

\subsection{Ghost} \label{sec:background_ghost}

Ghost~\cite{ghost_cms} is a modern, open-source publishing platform focused on independent journalism
and newsletters. It offers a clean editorial interface and native support for
membership models. Ghost defines four built-in roles --- \textit{Contributor},
\textit{Author}, \textit{Editor}, and \textit{Administrator} --- with well-defined
publish permissions. Contributors cannot publish directly and must submit content for
review. However, Ghost does not support explicit content rejection: editors can only
publish or leave content unpublished, with no formal rejection state or feedback
mechanism for contributors. This makes it unsuitable for a newsroom where editorial
accountability and structured feedback are part of the workflow. Its theming system
also makes full frontend decoupling impractical.

\subsection{Summary} \label{sec:background_summary}

The three platforms differ significantly in how they address the core requirements
of Diário LX.

Regarding \textbf{editorial workflow}, WordPress requires plugins to approximate any
structured content lifecycle, while Ghost supports a basic submit-and-publish model
where Contributors cannot publish directly but editors have no formal way to reject
content --- they can only publish or leave it unpublished, with no feedback mechanism
for contributors. Diário LX enforces a formal lifecycle with four explicit states ---
\textit{Draft}, \textit{Pending}, \textit{Published}, and \textit{Rejected} ---
requiring an editor to explicitly approve or reject each submission before publication.

Regarding \textbf{role-based access control}, WordPress relies on plugins for any
granular permission model. Ghost defines four built-in roles with well-defined publish
restrictions, but the permissions are tied to the platform's own publishing model and
cannot be extended. Diário LX defines three roles --- \textit{Contributor},
\textit{Editor}, and \textit{Administrator} --- enforced directly at the API level,
meaning every endpoint independently validates the caller's role without depending on
any plugin or external system.

Regarding \textbf{multimedia support}, both WordPress and Ghost treat multimedia as
a secondary concern, relying on basic media libraries or plugins. Diário LX includes
a native upload pipeline supporting images, video, and audio, with direct integration
into the content authoring workflow.

Regarding \textbf{frontend decoupling}, WordPress and Ghost both tie the presentation
layer to their own theming systems, which constrains the design and technology choices
available to the frontend. Diário LX exposes a dedicated REST API consumed by an
independent React application, meaning the frontend can be developed, deployed, and
replaced entirely without touching the backend.

Table~\ref{tab:cms_comparison} consolidates these differences.

\begin{table}[H]
      \centering
      \caption{Comparison of existing CMS solutions against Diário LX requirements.}
      \label{tab:cms_comparison}
      \begin{tabularx}{\textwidth}{lXXX}
            \toprule
            \textbf{Feature}   & \textbf{WordPress} & \textbf{Ghost}        & \textbf{Diário LX}      \\
            \midrule
            Editorial workflow & Plugin-based       & Submit/publish only,  & Draft → Pending →       \\
                               &                    & no rejection state    & Published / Rejected    \\
            \addlinespace
            Role-based access  & Plugin-based       & Four built-in roles,  & Three roles, enforced   \\
                               &                    & no rejection feedback & natively at API level   \\
            \addlinespace
            Multimedia support & Plugin-based       & Basic                 & Native (images, video,  \\
                               &                    &                       & audio)                  \\
            \addlinespace
            API-first design   & Partial            & Yes                   & Yes                     \\
            \addlinespace
            Custom frontend    & Theme-based        & Theme-based           & React (fully decoupled) \\
            \addlinespace
            Open source        & Yes                & Yes                   & Yes                     \\
            \bottomrule
      \end{tabularx}
\end{table}

None of the analysed off-the-shelf solutions fully satisfies the editorial workflow,
access control, and multimedia management requirements of Diário LX, motivating the
development of a purpose-built platform.